% document formatting
\documentclass[10pt]{article}
\usepackage[utf8]{inputenc}
\usepackage[left=1in,right=1in,top=1in,bottom=1in]{geometry}
\usepackage[T1]{fontenc}
\usepackage{xcolor}

% math symbols, etc.
\usepackage{amsmath, amsfonts, amssymb, amsthm}

% lists
\usepackage{enumerate}

% images
\usepackage{graphicx} % for images
\usepackage{tikz}

% code blocks
% \usepackage{minted, listings} 

% verbatim greek
\usepackage{alphabeta}

\graphicspath{{./assets/images/Week 1}}

\title{02-719 Week 1 \\ \large{Genomics and Epigenetics of the Brain}}
 
\author{Aidan Jan}

\date{\today}

\begin{document}
\maketitle

\section*{Logistics}
\begin{itemize}
	\item Office hours: Wed. 1-2 PM, GHC 7711
	\item Email: \texttt{apfenning@cmu.edu}
\end{itemize}

\section*{Motivation of Neurogenomics}
\textbf{Alzheimer's Disease}
\begin{itemize}
	\item More than 7 million Americans are currently living with Alzhiemer's.  (Projected to increase to 13 Million by 2050)
	\item 1 in 3 older adults dies with Alzheimer's or another dementia.  It kills more than breast and prostate cancer combined.
	\item \$409 Billion dollars on Alzheimer's spendings in 2026, projected to be \$1 Trillion in 2050
	\item The lifetime risk for Alzheimer's at age 45 for women is 1 in 5, and for men it is 1 in 10.  (Unknown reason)
\end{itemize}
\textbf{Parkinson's Disease}
\begin{itemize}
	\item Lower prevalence than Alzheimer's (1 in 100 over age 60)
	\item 1M people in the United States, and 5M people worldwide affected
\end{itemize}
\textbf{Addiction}
\begin{itemize}
	\item 46.3M people in the US have a substance use disorder in 2021.
	\item 105k drug overdose deaths per year
\end{itemize}
\textbf{Major Depressive Disorder}
\begin{itemize}
	\item 13.8M - 15.4M adults affected
	\item One of the leading causes of suicide
	\item Costs employers\$78.7 Billion per year
\end{itemize}

\subsection*{Disease Categories}
Most of these diseases fall into one of these following categories:
\begin{itemize}
	\item No cure
	\item No effective long-term treatment (Alzheimer's, Parkinson's, other neurodegenerative diseases)
	\item Treatment comes with severe side effects (Bipolar, Major Depressive Disorder, ADHD, other psychiatric disorders)
\end{itemize}

\subsection*{Why are Brain Disorders So Hard to Treat?}
The primary issue: What do you target when you treat a brain disorder?
\begin{itemize}
	\item Target a specific brain region
	\item Specific cell types
	\item Specific hormones and neurotransmitters
	\item Specific proteins
	\item Specific section of DNA (gene therapy)
	\item Specific neural pathway
\end{itemize}
The problem is that modern medicines for brain diseases are like sledge hammers.  One protein might influence the disease, but also regulate many other brain functions.  Same goes with neurotransmitters, genes, etc.  There isn't a single gene, protein, brain region, or biochemical pathway that equates to a single behavior.

\section*{Cellular Diversity of the Brain}
The cellular diversity of the brain is overwhelming.
\begin{itemize}
	\item Purkinje
	\item Microglia
	\item Astrocytes
	\item Bipolar
	\item Oligodendrocytes
	\item Schwann
\end{itemize}

\section*{Brief History of Neuroscience}
\begin{itemize}
	\item Santiago Ramon y Cajal (1852 - 1934): Traced neuron morphology and noticed there were different sub-types.  Nobel prize winner.
	\item Cajal had a modern Zeiss microscope, which scientific discovery would not be possible without.
	\item Cajal's work was more ``data-driven'' than ``hypothesis-driven''.  Rather than designing an experiment to test a specific theory of the brain, Cajal sought to characterize the diversity of neural cell types
	\item His work led to the neuron doctrine: ``The brain is made of many different types of neurons.''
\end{itemize}

\subsection*{Distinguishing Subtypes of Neurons}
\begin{itemize}
	\item Shape (Cytoarchitecture)
	\item Firing (Electrophysiology)
	\item Neurotransmitter (Glutamatergic vs. Gabaergic)
	\item Marker Genes (Somatostatin)
	\item Genomics
\end{itemize}
Single Neuron RNA-Seq measures genome-wide gene expression in single cells.

\subsubsection*{Epigenetics}
There is an extraordinary diversity of neural cells within the brain that have different molecular and physiological properties.  Differences in gene expression are caused by \textbf{epigenetic differences} in enhancers across neural cell types.  
\begin{itemize}
	\item There are differences in the genome sequence of enhancers that influence the evolution of complex traits (and human disease).
	\item Typical computational methods do not capture cell type-specific evolution at enhancer sequences.
	\item Machine learning has the potential to trace the evolution of neural cell type-specific enhancers across evolution.
\end{itemize}

\section*{Studying the Brain}
There are many challenges of studying the brain posed by cell types.
\begin{itemize}
	\item Deep brain stimulation / optogenetics - How do you selectively stimulate the desired cell type in a given brain region
	\item Neurological / Psychiatric disorders - What cell type(s) are underlying the disorder?  How do you target them?
	\item How do distinct subtypes of neurons form neural circuits for learning/memory?
\end{itemize}

\subsection*{Different Views of the Genome}
Computer scientists see the genome as a sequence of A's, T's, C's, and G's, and focus on identifying specific DNA sequences that might correlate with diseases. However, this view often does not capture the true complexity of cell types since it dismisses epigenetics completely.  Biologists not only see the genome, but also consider which parts might code for proteins, which might be transcriptionally active, which might be repressed.  Epigenetics, enhancers, and other regulatory elements are considered.

\subsection*{Epigenetic Mechanisms of Gene Regulation}
\begin{center} 
	\includegraphics*[width=0.7\textwidth]{W1_1.png} 
\end{center}
Different cell types might have the same gene, but regulated by different enhancers.  A given gene might have anywhere between 2 to over 100 enhancers regulating it.
\begin{itemize}
	\item The combination and spacing of \textbf{transcription factor binding sites} determines activation.
	\item Enhancers have a `grammer' to be able to operate correctly.
	\begin{itemize}
	    \item Their spacing, orientation, and order all matter.
	    \item Proteins bind to enhancer sequences on the DNA.  Two enhancers too close together or facing the wrong direction may cause the binding of one protein to inhibit the binding of the other protein.
    \end{itemize}
\end{itemize}
\begin{center} 
	\includegraphics*[width=0.5\textwidth]{W1_2.png} 
\end{center}
\begin{itemize}
	\item There are hundreds of thousands of enhancers across the human brain (and millions across the body)
	\item Most of these enhancers are cell type-specific
	\item Promoter regions and the genes themselves are much less specific
\end{itemize}

\subsection*{Conservation through Evolution}
Conservation methods are protein-centric.
\begin{itemize}
	\item If something has been maintained over millions of years of evolution, then there must be evolutionary constraint selecting against mutations in these regions.
	\item In protein-coding regions, every third base pair has a higher rate of mutating than the first two due to the \textbf{third codon wobble}; when observing the codons for different amino acids, the identity of the amino acid depends less on the third base pair compared to the first two.
	\item Conservation can be measured by PhyloP / PhastCons.  (Pollard et al., 2010)
\end{itemize}
Below is an example of how an enhancer might be linked to a physical trait.  
\begin{center} 
	\includegraphics*[width=\textwidth]{W1_3.png} 
\end{center}
\begin{itemize}
	\item The ZRS enhancer was discovered to be related to limb development when observing different organisms.  Snakes have the enhancer sequence deleted.
	\item When scientists deleted the same sequence from mice DNA, they ``Serpentized'' the mice, causing their limbs to not develop at birth.
\end{itemize}

\subsection*{Alignment of Individual Nucleotides is not Conservation}
Is Nucleotide-level conservation the exception or the rule?
\begin{center} 
	\includegraphics*[width=0.9\textwidth]{W1_4.png} 
\end{center}
\begin{itemize}
	\item It turns out that nucleotide-level conservation is \textbf{not} the rule.  Genes are selected for, not base pairs.
	\item In the different species on the diagram, they all share a Scaper enhancer which enhances a gene called Islet.  It turns out that this combination is present in all the species listed, and their nucleotide motifs are conserved, albeit appearing in different orders.
\end{itemize}
Nucleotide alignment does not capture regulatory evolution.
\begin{itemize}
	\item There is no strong difference in nucleotide-level evolution between enhancers that are conserved and divergent.
\end{itemize}

\section*{Machine Learning for Enhancer Prediction}
Studies have been done on training machine models on Motif scores, k-Mers, and raw sequences of nucleotides.  It turns out that raw sequence has by far the best performance.
\begin{itemize}
	\item Alternatively, use tissue-specific regulatory code.
\end{itemize}

\subsection*{How do we find Conserved Segments of the Genome?}
\begin{itemize}
	\item Convolutional neural network (deep learning) has revolutionized image analysis.
	\item We can adapt this method to make predictions of genome sequence function.
\end{itemize}
\begin{center} 
	\includegraphics*[width=\textwidth]{W1_5.png} 
\end{center}
Alternatively, use tissue-specific regulatory code
\begin{itemize}
	\item We can measure activity correlated with certain enhancer sequences in some species\dots
	\item and use it to predict activity in orthologous regions of other species.
\end{itemize}

\subsubsection*{The Zoonomia Consortium}
\begin{itemize}
	\item A series of papers sequencing DNA in brains (motor cortex) from multiple species, and predicting activity in motor cortex enhancers
\end{itemize}
\begin{center} 
	\includegraphics*[width=0.9\textwidth]{W1_6.png} 
\end{center}
\begin{itemize}
	\item The Y-axis are the open chromatin regions.
	\item The X-axis are different species.
	\item The input is the top graph, where only a few species are sampled for enhancer activity.
	\item The output is the bottom graph, essentially extrapolating from the input species.
	\item The CNN models are accurate, with >80\% validation accuracy.
\end{itemize}
The above project is called TACIT-ML (Tissue Aware Conservation Inference Through Machine Learning), and it outperforms previous models, such as PhyloP and PhastCons.
\begin{itemize}
	\item TACIT (predicted open chromatin) does a better job at identifying which mutations are associated with human disease and what cells they function in.
	\item TACIT can be used to identify part of the genome associated with the evolution of complex traits
\end{itemize}

\subsection*{Brain Size has Convergently Evolved}
There is an incredible diversity of mammalian brains.
\begin{itemize}
	\item It turns out that when plotted on a logarithmic scale graph (both axes), body correlates linearly with brain weight
\end{itemize}



\end{document}